\documentclass[11pt,a4paper]{article}

\usepackage[utf8]{inputenc}
\usepackage[T1]{fontenc}
\usepackage[french]{babel}
\usepackage{amsmath,amssymb,amsfonts}
\usepackage{geometry}
\usepackage{fancyhdr}
\usepackage{xcolor}
\usepackage{tcolorbox}
\usepackage{enumitem}
\usepackage{booktabs}
\usepackage{titlesec}
\usepackage{mdframed}
\usepackage{hyperref}

\geometry{top=1.8cm, bottom=1.8cm, left=2cm, right=2cm}

\definecolor{darkblue}{RGB}{31,56,100}
\definecolor{medblue}{RGB}{46,117,182}
\definecolor{lightblue}{RGB}{220,235,250}
\definecolor{darkred}{RGB}{160,30,30}
\definecolor{lightred}{RGB}{255,230,230}
\definecolor{darkgreen}{RGB}{30,120,50}
\definecolor{lightgreen}{RGB}{230,250,230}
\definecolor{lightyellow}{RGB}{255,250,220}
\definecolor{lightorange}{RGB}{255,240,220}

\tcbuselibrary{skins,breakable}

% Colored boxes
\newtcolorbox{formulabox}[1][]{colback=lightblue,colframe=medblue,fonttitle=\bfseries,title=#1,breakable}
\newtcolorbox{warningbox}[1][]{colback=lightred,colframe=darkred,fonttitle=\bfseries,title=#1,breakable}
\newtcolorbox{tipbox}[1][]{colback=lightgreen,colframe=darkgreen,fonttitle=\bfseries,title=#1,breakable}
\newtcolorbox{membox}[1][]{colback=lightyellow,colframe=orange!80!black,fonttitle=\bfseries,title=#1,breakable}

\pagestyle{fancy}
\fancyhf{}
\fancyhead[L]{\small\textcolor{gray}{Guide de R\'evision -- MOR 2025-2026}}
\fancyhead[R]{\small\textcolor{gray}{Examen 10 mars 2026}}
\fancyfoot[C]{\thepage}

\titleformat{\section}{\Large\bfseries\color{darkblue}}{\thesection}{1em}{}[\vspace{-0.3em}\color{medblue}\hrule\vspace{0.3em}]
\titleformat{\subsection}{\large\bfseries\color{medblue}}{\thesubsection}{1em}{}

\setlength{\parindent}{0pt}
\setlength{\parskip}{0.5em}

\begin{document}

% ============ TITLE ============
\begin{center}
{\Huge\bfseries\textcolor{darkblue}{Guide de R\'evision}}\\[0.3cm]
{\LARGE\textcolor{darkblue}{R\'eduction de Mod\`eles -- Examen}}\\[0.5cm]
{\large Mardi 10 mars 2026, 13h30, 2h, sans documents}\\[0.3cm]
{\large Masters (MS)$^2$SC et CHPS -- ENS Paris-Saclay}
\end{center}
\vspace{0.5cm}

% ============ PLANNING ============
\section{Plan de r\'evision optimis\'e}

\begin{center}
\begin{tabular}{@{}lll@{}}
\toprule
\textbf{Horaire} & \textbf{Bloc} & \textbf{Contenu} \\
\midrule
20h--21h30 & Bloc 1 & Lire ce guide entier (POD/SVD + PGD champ + PGD EDP) \\
21h30--23h & Bloc 2 & Refaire chaque question de l'examen 2025 \`a la main \\
23h--00h & Bloc 3 & Questions flash + formules cl\'es \\
8h--10h & Bloc 4 & Red\'eriver PGD (Q6, Q7, Q9) sur papier blanc \\
10h--11h30 & Bloc 5 & Relire formules, questions flash \\
11h30--12h & Bloc 6 & Dernier passage sur les pi\`eges courants \\
\bottomrule
\end{tabular}
\end{center}

% ============ PARTIE A ============
\section{Partie A : POD / SVD (Exercice 1 de l'examen)}

% === Q1 ===
\subsection{Q1 -- Interpr\'etation + d\'emonstration $S=XX^T$ + POD de rang $n$}

\textbf{\'Enonc\'e :} On a $N$ vecteurs $u^1,\ldots,u^N \in \mathbb{R}^m$ avec $m \gg N$. On note $X = \text{col}(u^i) \in \mathbb{R}^{m\times N}$.
\begin{enumerate}
\item Interpr\'eter $\displaystyle\min_{\|\varphi\|=1} \sum_{i=1}^{N} \|u^i - \langle u^i,\varphi\rangle\varphi\|^2$
\item Montrer que cela revient \`a \'etudier $S = XX^T$
\item \'Ecrire le probl\`eme POD de rang $n$
\end{enumerate}

\begin{formulabox}[Solution compl\`ete]
\textbf{Interpr\'etation :} On cherche la direction unitaire $\varphi$ qui minimise la somme des erreurs de projection des $u^i$ sur $\varphi$. C'est l'\textbf{axe principal} du nuage de points (analogue \`a l'ACP).

\medskip
\textbf{D\'emonstration :} Par le th\'eor\`eme de Pythagore~:
\[
\|u^i\|^2 = \|\langle u^i,\varphi\rangle\varphi\|^2 + \|u^i - \langle u^i,\varphi\rangle\varphi\|^2
\]
La norme $\|u^i\|^2$ \'etant constante, minimiser l'erreur $\Leftrightarrow$ maximiser la projection~:
\[
\max_{\|\varphi\|=1} \sum_{i=1}^{N} \langle u^i,\varphi\rangle^2
\]
Or $\langle u^i,\varphi\rangle = (u^i)^T\varphi$. Le vecteur $X^T\varphi \in \mathbb{R}^N$ regroupe ces produits scalaires~:
\[
\sum_{i=1}^{N} \big((u^i)^T\varphi\big)^2 = \|X^T\varphi\|^2 = (X^T\varphi)^T(X^T\varphi) = \varphi^T \underbrace{XX^T}_{S}\, \varphi
\]
C'est le \textbf{quotient de Rayleigh}~: $\displaystyle\max_{\|\varphi\|=1} \varphi^T S\varphi$. La solution est le vecteur propre de $S$ associ\'e \`a la plus grande valeur propre.

\medskip
\textbf{Propri\'et\'es de $S = XX^T$~:}
\begin{itemize}[nosep]
\item Dimension~: $m \times m$
\item Sym\'etrique~: $S^T = (XX^T)^T = XX^T = S$
\item Semi-d\'efinie positive~: $\forall v,\; v^TSv = \|X^Tv\|^2 \geq 0$
\end{itemize}

\medskip
\textbf{POD de rang $n$~:} On cherche $\Phi = [\varphi_1,\ldots,\varphi_n] \in \mathbb{R}^{m\times n}$, base orthonorm\'ee~:
\[
\boxed{\min_{\Phi^T\Phi = I_n} \sum_{i=1}^{N} \|u^i - \Phi\Phi^T u^i\|^2}
\]
Les $\varphi_k$ sont les $n$ vecteurs propres de $S$ associ\'es aux $n$ plus grandes valeurs propres.
\end{formulabox}

\begin{warningbox}[Pi\`eges Q1]
\begin{itemize}[nosep]
\item Ne pas oublier $\|\varphi\|=1$ (contrainte unitaire)
\item $S$ est $m\times m$, \textbf{PAS} $N\times N$
\item L'\'equivalence min erreur $\Leftrightarrow$ max projection passe par Pythagore
\end{itemize}
\end{warningbox}

% === Q2 ===
\subsection{Q2 -- Th\'eor\`eme SVD + SVD tronqu\'ee}

\begin{formulabox}[Solution]
\textbf{Th\'eor\`eme SVD~:} Toute matrice $X \in \mathbb{R}^{m\times N}$ admet~:
\[
\boxed{X = U\Sigma V^T}
\]
o\`u~:
\begin{itemize}[nosep]
\item $U \in \mathbb{R}^{m\times m}$ orthogonale ($U^TU = I_m$)~: vecteurs singuliers \`a gauche
\item $V \in \mathbb{R}^{N\times N}$ orthogonale ($V^TV = I_N$)~: vecteurs singuliers \`a droite
\item $\Sigma \in \mathbb{R}^{m\times N}$ pseudo-diagonale~: $\sigma_1 \geq \sigma_2 \geq \cdots \geq 0$
\end{itemize}

\textbf{SVD tronqu\'ee d'ordre $n$~:}
\[
X_n = U_n\Sigma_n V_n^T
\]
C'est la \textbf{meilleure approximation de rang $n$} au sens de Frobenius (th. d'Eckart--Young)~:
\[
\|X - X_n\|_F = \sqrt{\sum_{k=n+1}^{r} \sigma_k^2}
\]
\end{formulabox}

% === Q3 ===
\subsection{Q3 -- Lien POD $=$ SVD}

\begin{formulabox}[Solution -- D\'emonstration cl\'e]
On injecte $X = U\Sigma V^T$ dans $S$~:
\[
S = XX^T = (U\Sigma V^T)(V\Sigma^T U^T) = U(\Sigma\Sigma^T)U^T = U\Lambda U^T
\]
car $V^TV = I_N$.

\medskip
\textbf{Conclusions~:}
\begin{enumerate}[nosep]
\item Colonnes de $U$ (vect.\ singuliers gauches) $=$ vecteurs propres de $S$ (modes POD)
\item $\boxed{\lambda_i = \sigma_i^2}$ \quad (valeurs propres de $S$ $=$ carr\'es des valeurs singuli\`eres)
\end{enumerate}

\textbf{Calcul pratique (SVD via $S$)~:}
\begin{enumerate}[nosep]
\item Former $S = XX^T$
\item Diagonaliser $S \Rightarrow U$ et $\Lambda$
\item $\sigma_i = \sqrt{\lambda_i}$
\item $v_i = \dfrac{1}{\sigma_i}X^Tu_i$ \quad (reconstituer $V$ colonne par colonne)
\end{enumerate}
\end{formulabox}

% === Q4 ===
\subsection{Q4 -- Utiliser \texttt{svd()} pour l'approximation \`a variables s\'epar\'ees}

\begin{formulabox}[Solution]
Soit $u(t,x)$ d\'efini sur $[0,T]\times[0,L]$, discr\'etis\'e en $N$ instants et $m$ points.

\begin{enumerate}[nosep]
\item Construire la matrice des snapshots $X \in \mathbb{R}^{m\times N}$ avec $X_{i,j} = u(x_i,t_j)$
\item Calculer \texttt{[U,S,V] = svd(X)}
\item Approximation s\'epar\'ee~:
\[
\boxed{u(x,t) \approx \sum_{k=1}^{n} \sigma_k\, U_k(x)\, V_k(t)}
\]
\item Choix de $n$ (crit\`ere d'\'energie)~: le plus petit $n$ tel que~:
\[
\boxed{1 - \frac{\sum_{k=1}^{n}\sigma_k^2}{\sum_{k=1}^{N}\sigma_k^2} \leq \varepsilon}
\]
\end{enumerate}
\end{formulabox}

% === Q5 ===
\subsection{Q5 -- $n$-\'epaisseur de Kolmogorov}

\begin{formulabox}[Solution]
\textbf{D\'efinition formelle~:}
\[
\boxed{d_n(\mathcal{U},V) = \inf_{\substack{V_n \subset V \\ \dim(V_n)=n}} \sup_{u \in \mathcal{U}} \inf_{v \in V_n} \|u-v\|_V}
\]
C'est la pire erreur d'approximation pour le meilleur sous-espace de dimension $n$.

\medskip
\textbf{Algorithme pragmatique~:}
\begin{enumerate}[nosep]
\item \'Echantillonnage~: g\'en\'erer $M$ solutions $\Rightarrow$ matrice $X$
\item SVD~: $X = U\Sigma V^T$
\item Estimation~: $\boxed{d_n(X) = \|X - X_n\|_2 = \sigma_{n+1}}$
\end{enumerate}
\end{formulabox}

\begin{warningbox}[Pi\`ege]
C'est $\sigma_{n+1}$ (pas $\sigma_n$) car c'est l'erreur \textbf{apr\`es} avoir gard\'e $n$ modes.
\end{warningbox}

% ============ PARTIE B ============
\section{Partie B : PGD pour un champ donn\'e (Exercice 2)}

\subsection{Q6 -- Construction de l'approximation PGD de $u(x,t)$}

\textbf{\'Enonc\'e :} Construire $u_n(x,t) = \sum_{l=1}^{n}\Lambda_l(x)\lambda_l(t)$ par algorithme glouton.

\begin{formulabox}[Solution compl\`ete]
\textbf{Principe glouton~:} On conna\^{\i}t $u_{n-1}$. On pose $e = u - u_{n-1}$ (r\'esidu). On cherche $R(x)S(t)$ minimisant~:
\[
J(R,S) = \frac{1}{2}\int_I\int_\Omega \big(e(x,t) - R(x)S(t)\big)^2\,\mathrm{d}x\,\mathrm{d}t
\]

\textbf{R\'esolution par directions altern\'ees~:}

\medskip
\underline{\'Etape 1~: $S(t)$ fix\'e, trouver $R(x)$}

Variation par rapport \`a $R$ (perturbation $R^*$)~:
\[
\int_I\int_\Omega (e - RS)\,R^*S\,\mathrm{d}x\,\mathrm{d}t = 0 \quad\forall R^*
\]
On factorise $R^*$ et on sort les int\'egrales temporelles~:
\[
\int_\Omega R^*\left[\int_I eS\,\mathrm{d}t - R\int_I S^2\,\mathrm{d}t\right]\mathrm{d}x = 0 \quad\forall R^*
\]
Par le lemme fondamental du calcul des variations~:
\[
\boxed{R(x) = \frac{\displaystyle\int_I\big(u(x,t)-u_{n-1}(x,t)\big)S(t)\,\mathrm{d}t}{\displaystyle\int_I S^2(t)\,\mathrm{d}t}}
\]

\underline{\'Etape 2~: $R(x)$ fix\'e, trouver $S(t)$}

Par sym\'etrie~:
\[
\boxed{S(t) = \frac{\displaystyle\int_\Omega\big(u(x,t)-u_{n-1}(x,t)\big)R(x)\,\mathrm{d}x}{\displaystyle\int_\Omega R^2(x)\,\mathrm{d}x}}
\]

\textbf{Algorithme~:}
\begin{enumerate}[nosep]
\item $\Delta U = u - u_{n-1}$
\item Initialiser $S(t)$ (ex: $S(t) = t$), poser $s=100$
\item Tant que $s > 10^{-3}$~: (a) $S_\text{old} = S_\text{new}$ \quad (b) Calculer $R(x)$, normaliser \quad (c) Calculer $S(t)$ \quad (d) $s = \frac{\int_I(S_\text{new}-S_\text{old})^2\,\mathrm{d}t}{\int_I S_\text{old}^2\,\mathrm{d}t}$
\item Stocker $(\Lambda_i, \Gamma_i) = (R,S)$
\end{enumerate}
\end{formulabox}

\begin{warningbox}[Pi\`eges Q6]
\begin{itemize}[nosep]
\item \textbf{Normaliser} $R$ apr\`es chaque calcul (convention~: normaliser le spatial)
\item Le d\'enominateur est un \textbf{scalaire} (norme $L^2$ au carr\'e de la fonction fix\'ee)
\item Ne pas oublier l'orthogonalisation de \textbf{Gram--Schmidt} apr\`es chaque mode
\end{itemize}
\end{warningbox}

% === Q7 ===
\subsection{Q7 -- PGD pour Poisson $-\Delta u = f(x,t)$}

\begin{formulabox}[Solution compl\`ete]
\textbf{Formulation variationnelle~:} Multiplier par $u^*$ (nulle sur $\partial\Omega$), int\'egrer, Green~:
\[
\int_I\int_\Omega \nabla u \cdot \nabla u^*\,\mathrm{d}x\,\mathrm{d}t = \int_I\int_\Omega f\,u^*\,\mathrm{d}x\,\mathrm{d}t
\]

\textbf{PGD~:} $u = u_{n-1} + R(x)S(t)$, \quad $u^* = R^*(x)S(t) + R(x)S^*(t)$

\medskip
\underline{\'Etape spatiale ($S$ fix\'e, $u^* = R^*S$)~:}
\[
\underbrace{\left(\int_I S^2\,\mathrm{d}t\right)}_{\text{scalaire}} \int_\Omega\nabla R\cdot\nabla R^*\,\mathrm{d}x = \int_\Omega R^*\underbrace{\left(\int_I fS\,\mathrm{d}t\right)}_{\text{fct de }x}\mathrm{d}x - \int_\Omega\nabla R^*\cdot\underbrace{\left(\int_I\nabla u_{n-1}\,S\,\mathrm{d}t\right)}_{\text{fct de }x}\mathrm{d}x
\]
$\Rightarrow$ C'est un \textbf{probl\`eme de Poisson standard} pour $R(x)$ !

\medskip
\underline{\'Etape param\'etrique ($R$ fix\'e, $u^* = RS^*$)~:}
\[
\boxed{S(t)\int_\Omega|\nabla R|^2\,\mathrm{d}x = \int_\Omega fR\,\mathrm{d}x - \int_\Omega\nabla u_{n-1}\cdot\nabla R\,\mathrm{d}x}
\]
$\Rightarrow$ C'est une \textbf{\'equation alg\'ebrique} pour $S(t)$ ! (pas d'ODE car pas de d\'eriv\'ee temporelle)

\medskip
\textbf{Acc\'el\'eration (Update)~:} Apr\`es avoir calcul\'e $\{\Lambda_1,\ldots,\Lambda_n\}$, on fige la base spatiale et on recalcule \textbf{simultan\'ement} toutes les $\lambda_1(t),\ldots,\lambda_n(t)$ via un syst\`eme coupl\'e $n\times n$.
\end{formulabox}

% === Q8 ===
\subsection{Q8 -- Conditions aux limites non nulles}

\begin{formulabox}[Solution]
PGD cherche $u = \sum\Lambda_l(x)\lambda_l(t)$. Si $u=g\neq 0$ sur le bord, il faudrait une \textbf{infinit\'e de modes} pour reproduire $g$ par produit tensoriel.

\textbf{Traitement par rel\`evement~:}
\begin{enumerate}[nosep]
\item D\'efinir $u_g(x,t)$ satisfaisant les CL
\item Poser $u = \tilde{u} + u_g$
\item $\tilde{u}$ v\'erifie des CL homog\`enes ($\tilde{u}=0$ sur $\partial\Omega$)
\item \'Equation modifi\'ee~: $-\Delta\tilde{u} = f + \Delta u_g$
\end{enumerate}
Pour la chaleur~: $M\dot{W} + KW = F - M\dot{u}_\text{CL} - Ku_\text{CL}$
\end{formulabox}

% ============ PARTIE C ============
\section{Partie C : PGD pour une EDP parabolique (Exercice 3)}

\subsection{Q9 -- PGD pour $\partial_{tt}u + \partial_x u + u = r(x,t)$}

\begin{formulabox}[Solution compl\`ete]
\textbf{\'Equation~:} $\partial_{tt}u + \partial_x u + u = r(x,t)$ sur $[0,1]\times[0,T]$, CL et CI nulles.

\medskip
\textbf{\'Etape 1 -- Formulation variationnelle~:}

On cherche $u(x,t) \approx X(x)T(t)$, fonction test $u^* = X^*T + XT^*$~:
\[
\int_0^T\!\!\int_0^1 \big(XT'' + X'T + XT\big)\big(X^*T + XT^*\big)\,\mathrm{d}x\,\mathrm{d}t = \int_0^T\!\!\int_0^1 r\big(X^*T + XT^*\big)\,\mathrm{d}x\,\mathrm{d}t
\]

\textbf{\'Etape 2 -- Sous-probl\`eme spatial} ($T$ connu, test $X^*T$)~:

On identifie terme par terme~:
\begin{itemize}[nosep]
\item $XT''$ test\'e par $X^*T$ $\Rightarrow$ $\displaystyle\underbrace{\int_0^T T''T\,\mathrm{d}t}_{\alpha_1}\int_0^1 XX^*\,\mathrm{d}x$
\item $X'T$ test\'e par $X^*T$ $\Rightarrow$ $\displaystyle\underbrace{\int_0^T T^2\,\mathrm{d}t}_{\beta}\int_0^1 X'X^*\,\mathrm{d}x$
\item $XT$ test\'e par $X^*T$ $\Rightarrow$ $\displaystyle\beta\int_0^1 XX^*\,\mathrm{d}x$
\end{itemize}
Probl\`eme spatial~:
\[
\boxed{\int_0^1\big(\alpha_1 X + \beta X' + \beta X\big)X^*\,\mathrm{d}x = \int_0^1 \tilde{r}(x)\,X^*\,\mathrm{d}x}
\]
avec $\tilde{r}(x) = \int_0^T r(x,t)T(t)\,\mathrm{d}t$ $\Rightarrow$ \textbf{EDO spatiale 1D}

\medskip
\textbf{\'Etape 3 -- Sous-probl\`eme temporel} ($X$ connu, test $XT^*$)~:

\begin{itemize}[nosep]
\item $A = \int_0^1 X^2\,\mathrm{d}x$, \quad $B = \int_0^1(X'X + X^2)\,\mathrm{d}x$, \quad $\hat{r}(t) = \int_0^1 r(x,t)X\,\mathrm{d}x$
\end{itemize}
\[
\boxed{\int_0^T (AT'' + BT)T^*\,\mathrm{d}t = \int_0^T \hat{r}\,T^*\,\mathrm{d}t} \qquad\Rightarrow\qquad AT'' + BT = \hat{r}(t)
\]
$\Rightarrow$ \textbf{EDO temporelle d'ordre 2}

\medskip
\textbf{Co\^ut de calcul~:}
\begin{itemize}[nosep]
\item Classique~: matrice $(N_x\!\times\!N_t)^2$ $\Rightarrow$ co\^ut $\mathcal{O}((N_xN_t)^p)$ \textbf{(multiplicatif)}
\item PGD~: s\'epar\'ement $N_x\!\times\!N_x$ puis $N_t\!\times\!N_t$ $\Rightarrow$ co\^ut $\mathcal{O}(N_x^p + N_t^p)$ \textbf{(additif)}
\end{itemize}
La PGD \textbf{contourne la mal\'ediction de la dimensionnalit\'e}.
\end{formulabox}

% ============ PARTIE D ============
\section{Partie D : PGD pour l'\'equation de la chaleur (TP2)}

\begin{formulabox}[Sch\'ema complet -- $\rho c\,\partial_t u - k\,\partial_{xx}u = f(x,t)$]

\textbf{Forme semi-discr\`ete (FEM spatial)~:} $M\dot{u} + Ku = F$

\textbf{Changement de variable (CL non nulles)~:} $u = u_\text{CL} + W$
\[
M\dot{W} + KW = Q_0 = F - M\dot{u}_\text{CL} - Ku_\text{CL}
\]

\textbf{PGD~:} $W \approx \sum_i \Lambda_i(x)\lambda_i(t)$

\medskip
\underline{Sous-probl\`eme spatial ($\lambda$ fix\'e)~:}
\[
\boxed{J\,\Lambda = q}
\]
\[
J = \left[\int_0^T\lambda\dot{\lambda}\,\mathrm{d}t\right]M + \left[\int_0^T\lambda^2\,\mathrm{d}t\right]K \qquad q = \int_0^T Q(t)\lambda(t)\,\mathrm{d}t
\]
Syst\`eme FE de taille $N_x\times N_x$. Apr\`es r\'esolution, \textbf{normaliser} $\Lambda$.

\medskip
\underline{Sous-probl\`eme temporel ($\Lambda$ fix\'e)~:}
\[
\boxed{\alpha\,\dot{\lambda} + \beta\,\lambda = \gamma(t)}
\]
\[
\alpha = \Lambda^T M\Lambda \quad(\text{scalaire}), \qquad \beta = \Lambda^T K\Lambda \quad(\text{scalaire}), \qquad \gamma(t) = \Lambda^T Q(t)
\]
ODE scalaire ! R\'esolution par \textbf{Euler implicite}~:
\[
\boxed{\lambda_i = \frac{\gamma(t_i)\,\Delta t + \alpha\,\lambda_{i-1}}{\alpha + \beta\,\Delta t}}
\]
\end{formulabox}

% ============ FORMULES CLES ============
\section{Formules cl\'es \`a m\'emoriser}

\begin{membox}[POD / SVD]
\begin{align*}
S &= XX^T \quad\text{(matrice de covariance, }m\times m\text{)} \\
S &= U\Lambda U^T \quad\text{(diagonalisation = POD)} \\
X &= U\Sigma V^T \quad\text{(SVD)} \\
\lambda_i &= \sigma_i^2 \quad\text{(lien POD--SVD)} \\
v_i &= \tfrac{1}{\sigma_i}X^Tu_i \quad\text{(reconstituer }V\text{)} \\
\varepsilon(n) &= \sqrt{\frac{\sum_{k=n+1}^{r}\sigma_k^2}{\sum_{k=1}^{r}\sigma_k^2}} \quad\text{(erreur relative)} \\
d_n &\approx \sigma_{n+1} \quad\text{(Kolmogorov)}
\end{align*}
\end{membox}

\begin{membox}[PGD -- champ donn\'e]
\[
R(x) = \frac{\int_I\Delta U\cdot S\,\mathrm{d}t}{\int_I S^2\,\mathrm{d}t} \qquad\qquad S(t) = \frac{\int_\Omega\Delta U\cdot R\,\mathrm{d}x}{\int_\Omega R^2\,\mathrm{d}x}
\]
\end{membox}

\begin{membox}[PGD -- \'Eq.\ chaleur]
\begin{align*}
\text{Spatial~:} & \quad J\Lambda = q, \quad J = \Big[\textstyle\int\lambda\dot{\lambda}\,\mathrm{d}t\Big]M + \Big[\textstyle\int\lambda^2\,\mathrm{d}t\Big]K \\
\text{Temporel~:} & \quad \alpha\dot{\lambda}+\beta\lambda=\gamma, \quad \alpha=\Lambda^TM\Lambda,\; \beta=\Lambda^TK\Lambda,\; \gamma=\Lambda^TQ \\
\text{Euler impl.~:} & \quad \lambda_i = (\gamma_i\Delta t + \alpha\lambda_{i-1})/(\alpha+\beta\Delta t)
\end{align*}
\end{membox}

\begin{membox}[Matrices \'el\'ementaires FE 1D]
\[
\underline{\underline{m}} = \frac{\Delta x}{6}\begin{pmatrix}2&1\\1&2\end{pmatrix} \qquad \underline{\underline{k}} = \frac{1}{\Delta x}\begin{pmatrix}1&-1\\-1&1\end{pmatrix}
\]
\end{membox}

% ============ QUESTIONS FLASH ============
\section{Questions flash (30 secondes chacune)}

\begin{enumerate}[leftmargin=*]
\item \textbf{Dimension de $S=XX^T$ si $X$ est $m\times N$ ?} \hfill $m\times m$

\item \textbf{Pourquoi $S$ est SDP ?} \hfill $v^TSv = \|X^Tv\|^2 \geq 0$

\item \textbf{Lien valeurs propres de $S$ / valeurs singuli\`eres de $X$ ?} \hfill $\lambda_i = \sigma_i^2$

\item \textbf{Estimation de la $n$-\'epaisseur de Kolmogorov ?} \hfill $\sigma_{n+1}$

\item \textbf{D\'enominateur de la formule PGD pour $R(x)$ ?} \hfill $\int_I S^2\,\mathrm{d}t$ (scalaire)

\item \textbf{CL non nulles dans la PGD ?} \hfill Rel\`evement : $u = \tilde{u}+u_g$

\item \textbf{\'Etape temporelle PGD pour Poisson ?} \hfill \'Equation \textbf{alg\'ebrique}

\item \textbf{\'Etape temporelle PGD pour la chaleur ?} \hfill ODE~: $\alpha\dot{\lambda}+\beta\lambda=\gamma$

\item \textbf{Avantage en co\^ut de la PGD ?} \hfill Additif $\mathcal{O}(N_x^p+N_t^p)$ vs multiplicatif

\item \textbf{Qu'est-ce que l'Update ?} \hfill Figer base spatiale, recalculer toutes les $\lambda_i(t)$

\item \textbf{Normalisation de $\Lambda$ ?} \hfill $\Lambda \leftarrow \Lambda/\sqrt{\Lambda^T I_x\Lambda}$

\item \textbf{Crit\`ere de troncature SVD ?} \hfill $1-\frac{\sum_{k=1}^n\sigma_k^2}{\sum\sigma_k^2}\leq\varepsilon$

\item \textbf{Matrice $J$ du probl\`eme spatial (chaleur) ?} \hfill $[\int\lambda\dot{\lambda}\,\mathrm{d}t]\,M + [\int\lambda^2\,\mathrm{d}t]\,K$

\item \textbf{Erreur en norme $L^2(\Omega\times I)$ ?} \hfill Int\'egrer en temps, puis en espace

\item \textbf{\'Etape temporelle pour $\partial_{tt}u+\partial_x u+u=r$ ?} \hfill ODE d'ordre \textbf{2}~: $AT''+BT=\hat{r}$
\end{enumerate}

% ============ ERREURS CLASSIQUES ============
\section{Erreurs classiques \`a \'eviter}

\begin{warningbox}[10 pi\`eges fr\'equents]
\begin{enumerate}[nosep]
\item Confondre $S=XX^T$ ($m\!\times\!m$) avec $X^TX$ ($N\!\times\!N$)
\item Oublier de \textbf{normaliser} $\Lambda$ apr\`es le sous-probl\`eme spatial
\item Confondre l'\'etape temporelle~: Poisson $\to$ alg\'ebrique, chaleur $\to$ ODE ordre 1, $\partial_{tt}$ $\to$ ODE ordre 2
\item Oublier le \textbf{rel\`evement} quand les CL sont non nulles
\item Ne pas \'ecrire la fonction test s\'epar\'ee compl\`ete~: $u^* = R^*S + RS^*$ (les \textbf{deux} termes)
\item Oublier que le d\'enominateur des formules PGD est un \textbf{scalaire}
\item Kolmogorov~: c'est $\sigma_{n+1}$, pas $\sigma_n$
\item Co\^ut PGD~: \textbf{additif} (pas multiplicatif)
\item Gram--Schmidt~: quand on retire de l'info de $\Lambda_{m+1}$, il faut compenser dans $\Gamma_i$
\item Matrice $J$ (chaleur)~: le terme $\int\lambda\dot{\lambda}\,\mathrm{d}t$ multiplie $M$, pas $K$
\end{enumerate}
\end{warningbox}

% ============ STRATEGIE ============
\section{Strat\'egie d'examen}

\begin{tipbox}[Gestion du temps (2h)]
\begin{itemize}[nosep]
\item \textbf{5 min}~: Lire tout le sujet
\item \textbf{30--40 min}~: Exercice 1 (POD/SVD) -- points ``gratuits'', bien r\'ediger Pythagore + Rayleigh
\item \textbf{40 min}~: Exercice 2 (PGD champ/Poisson) -- d\'erivations par directions altern\'ees
\item \textbf{40 min}~: Exercice 3 (PGD EDP) -- le plus dur mais structure identique
\item \textbf{5 min}~: Relecture
\end{itemize}
\textbf{En cas de blocage~:} \'ecrire la formulation variationnelle + la forme s\'epar\'ee de $u$ et $u^*$. Le reste suit m\'ecaniquement.
\end{tipbox}

\vfill
\begin{center}
\Large\textbf{Bonne chance pour demain !}
\end{center}

\end{document}
